

\def\PUDcodigo{ACE005}

\if\COMPILAR0
    \def\PUDcodacad{04507.4}
\else
    \def\PUDcodacad{04506.4}
\fi

\def\PUDdisciplina{Química Geral}

\def\PUDsemestre{S1}

\def\PUDhoras{80}

%NOVOS ---------------------------------------------------
\def\PUDhorasTeoria{60}
\def\PUDhorasPratica{20}

\def\PUDinterECA{---}

\if\COMPILAR0
   \def\PUDinterECA{ \hyperref[PUD:EME108]{Materiais (04507.22)} }
\else
   \def\PUDinterECA{ \hyperref[PUD:EME108]{Materiais I (04506.20)} }
\fi

\def\PUDrecursos{Projetor multimídia; Quadro branco e pincel; Aulas em laboratório para realização de experimentos químicos.}

%----------------------------------------------------------

\def\PUDcreditos{4}

\def\PUDprerequisito{---}

\def\PUDpreacad{---}

\def\PUDobjetivos{Familiarizar-se à área de química através dos seguintes objetivos específicos: Realizar cálculos químicos em reações químicas e solução; Demostrar natureza da radiação eletromagnética, suas características e seu efeito sobre os metais; Construir um modelo de estrutura do átomo justificando as suas propriedades para cada elemento químico; Projetar a formação de substâncias químicas utilizando os modelos de ligação química; Construir modelos representativos dos estados sólido, líquido e gasoso conforme suas propriedades; Interpretar os processos químicos baseado na lei da conservação da energia (1ª Lei da Termodinâmica); Explicar as reações de oxi-redução que podem ser usadas para gerar eletricidade, obter metais e proteger materiais; Usar os conceitos de  ácido-base nos cálculos de pH e em reações de neutralização.}

\def\PUDementa{Estequiometria.  Natureza da luz.  Estrutura do átomo.  Ligações químicas.  Estados da matéria.  Termoquímica.  Eletroquímica.  Ácidos e Bases.}

\def\PUDconteudo{ &  Natureza da luz: Características da radiação eletromagnética.  Quanta e fótons.  O efeito fotoelétrico. Estrutura eletrônica do átomo: O espectro de linhas do átomo de hidrogênio e o modelo de Bohr.  A dualidade onda-partícula da matéria.  O princípio da incerteza.  Orbitais atômicos.  Energia dos orbitais e os espectros atômicos.  Orbitais moleculares.  Tamanhos atômicos.  Energia de ionização e afinidade eletrônica.
\\
 &  Ligações químicas: Ligações iônicas.  A formação de íons.  Energia de rede.  Ligações covalentes.  Descrição da ligação covalente.  Energia e comprimento da ligação.  Eletronegatividade e polaridade de ligações.  Ligações metálicas.  Teoria das bandas.  Isolantes, semicondutores e condutores.  Semicondutores dopados tipo n e tipo p.
\\
 &  Estequiometria: Escrevendo e balanceando as equações químicas.  Estequiometria de reações químicas.  Conceito de mol e massa molar.  Soluções e concentração de soluções.  Cálculos químicos.  Reagente limitante.  Rendimento percentual. 
\\ 
 &  Estados da matéria: Modelos cinéticos molecular dos sólidos, líquidos e gases.  Gases.  Pressão gasosa.  Leis dos gases.  Mistura de gases.  Forças intermoleculares.  Propriedades dos líquidos.  Tensão superficial.  Viscosidade.  Estrutura do sólido.  Classificação dos sólidos.  Células unitárias e difração de raios X.   Sólidos metálicos.  Sólidos iônicos.  Sólidos covalentes e moleculares.  Diagrama de fase. 
\\ 
 &  Termoquímica: Sistema, fronteira e vizinhança.  Processos endotérmicos e exotérmicos.  Função de estado.  Calorimetria.  Entalpia de reação.  Entalpia de Combustão e poder calorífico de um combustível. 
\\ 
 &  Eletroquímica: Oxidação e redução.  Números de oxidação.  Agentes oxidantes e redutores.  Meias reações.  Célula eletroquímica.  Potenciais padrão de eletrodo e potencial padrão de célula.  Pilhas e baterias.  Eletrodeposição.  Corrosão. 
\\ 
 & Ácidos e bases: Ácidos e bases em solução aquosa.  Ácidos e bases fortes e fracos.  Reação de neutralização.  Escala de PH. }

\def\PUDmetodo{Aulas expositórias que podem ser teóricas e/ou práticas, onde as práticas no laboratório serão marcadas no decorrer da disciplina.}

\def\PUDavalia{A avaliação será feita com aplicação de provas teóricas e/ou práticas, além da possibilidade de inclusão de trabalhos e seminários no decorrer da disciplina.}

\def\PUDbibbasica{ &  \href{www.biblioteca.ifce.edu.br/index.asp?codigo_sophia=37436}{Farias}, Robson Fernandes de.  Química Geral no Contexto das Engenharias.  Campinas - SP.  Editora ATOMO.  2011.  135p.  ISBN: 9788576701675.
%&  BROWN, L.  S.;  HOLME, T.  A.;  Química Geral Aplicada à Engenharia.  1ª edição.  São Paulo.  Editora CENGAGE.  2009.  655p.  ISBN: 9788522106882.
\\
 &  \href{www.biblioteca.ifce.edu.br/index.asp?codigo_sophia=24474}{Brown}, Theodore L.  Química: a ciência central.  9º ed.  São Paulo, SP: Pearson Prentice Hall, 2005.  972p.  ISBN: 9788587918420.
\\
 &  \href{www.biblioteca.ifce.edu.br/index.asp?codigo_sophia=23790}{Kotz}, John C.  Química geral e reações químicas.  São Paulo, SP : Cengage Learning, 2008.  671p.  ISBN: 8522104271. }

\def\PUDbibcomplementar{ &  \href{www.biblioteca.ifce.edu.br/index.asp?codigo_sophia=37868}{Ball}, David W.  Físico-química: volume 1.  São Paulo, SP : Cengage Learning, 2013.  450p.  ISBN: 9788522104178.
%&  Chang, Raymond.  Química geral : conceitos essenciais.  4ª ed.  Porto Alegre, RS : AMGH, 2010.  778p.  ISBN: 9788563308047. 
\\ 
 & \href{www.biblioteca.ifce.edu.br/index.asp?codigo_sophia=24432}{Brady}, James E.  Química geral: volume 1.  2º ed.  Rio de Janeiro, RJ : LTC, 2006.  410p.  ISBN: 8521604489. 
\\
 & \href{www.biblioteca.ifce.edu.br/index.asp?codigo_sophia=23687}{Masterton}, William L.  Princípios de Química.  6ª ed.  Rio de Janeiro - RJ.  Editora LTC.  1990.  681p.  ISBN: 9788521611219. 
\\
 &  \href{www.biblioteca.ifce.edu.br/index.asp?codigo_sophia=24299}{Russell}, John B.  Química geral - v. 1.  2º ed.  São Paulo, Pearson Makron Books, 1994.  ISBN: 8534601925.
\\  
 &  \href{www.biblioteca.ifce.edu.br/index.asp?codigo_sophia=61743}{Chang}, Raymond.  Química geral: conceitos essenciais.  4º ed.  Porto Alegre, RS: AMGH, 2010.  778p.  ISBN: 9788563308047. }
 
\def\PUDautor{Samuel Vieira Dias}

\def\PUDdata{2013-04-23}

\def\PUDrevisao{1}

\def\PUDrevisor{Samuel Vieira Dias}

\def\PUDdatarevisao{2017-05-09}

\PUDcorpo
